% ============================================================
%  Manger Consciemment — Guide pratique sur les additifs alimentaires
%  faustobe.it | 2026 | CC BY-NC 4.0
% ============================================================
\documentclass[a5paper,12pt,oneside,openany]{book}

\usepackage[T1]{fontenc}
\usepackage[utf8]{inputenc}
\usepackage{ebgaramond}
\usepackage[a5paper, top=2cm, bottom=3cm, inner=2.2cm, outer=3cm]{geometry}
\usepackage[table]{xcolor}
\usepackage[french]{babel}

% Colors
\definecolor{verdescuro}{HTML}{2d6a4f}
\definecolor{verdecharo}{HTML}{52b788}
\definecolor{grigioscuro}{HTML}{2b2d42}
\definecolor{jaune}{HTML}{d4a017}
\definecolor{rouge}{HTML}{c1121f}
\definecolor{bleuinfo}{HTML}{3a86ff}

\usepackage{tcolorbox}
\tcbuselibrary{skins,breakable}

\usepackage{booktabs}
\usepackage{tabularx}
\newcolumntype{L}{>{\raggedright\arraybackslash}X}

\usepackage{hyperref}
\hypersetup{
  colorlinks=true,
  linkcolor=verdescuro,
  urlcolor=verdescuro,
  pdftitle={Manger Consciemment},
  pdfauthor={faustobe},
}

\usepackage{qrcode}
\usepackage{fancyhdr}
\usepackage{titlesec}
\usepackage{microtype}

% Colored boxes
\tcbset{
  boitebase/.style={
    breakable, enhanced, arc=4pt, boxrule=0.8pt,
    fonttitle=\bfseries\small,
    before upper={\parskip=4pt},
  }
}

\newcommand{\boxvert}[2]{%
  \begin{tcolorbox}[boitebase,
    colback=verdescuro!8!white, colframe=verdescuro, title={#1}]
  #2\end{tcolorbox}}

\newcommand{\boxjaune}[2]{%
  \begin{tcolorbox}[boitebase,
    colback=jaune!12!white, colframe=jaune!80!black, title={#1}]
  #2\end{tcolorbox}}

\newcommand{\boxrouge}[2]{%
  \begin{tcolorbox}[boitebase,
    colback=rouge!8!white, colframe=rouge, title={#1}]
  #2\end{tcolorbox}}

\newcommand{\boxinfo}[2]{%
  \begin{tcolorbox}[boitebase,
    colback=bleuinfo!8!white, colframe=bleuinfo, title={#1}]
  #2\end{tcolorbox}}

% Page style
\pagestyle{fancy}
\fancyhf{}
\fancyfoot[C]{\footnotesize\textcolor{grigioscuro}{\textit{Manger Consciemment\ \textbar\ faustobe.it}}}
\fancyfoot[R]{\footnotesize\thepage}
\renewcommand{\headrulewidth}{0pt}
\renewcommand{\footrulewidth}{0.4pt}
\fancypagestyle{plain}{%
  \fancyhf{}%
  \fancyfoot[C]{\footnotesize\textcolor{grigioscuro}{\textit{Manger Consciemment\ \textbar\ faustobe.it}}}%
  \fancyfoot[R]{\footnotesize\thepage}%
  \renewcommand{\headrulewidth}{0pt}%
  \renewcommand{\footrulewidth}{0.4pt}%
}

% Chapter title style
\titleformat{\chapter}[display]
  {\normalfont\LARGE\bfseries\color{verdescuro}}
  {\chaptertitlename\ \thechapter}{10pt}{\Huge}
\titlespacing*{\chapter}{0pt}{-20pt}{30pt}

% ============================================================
\begin{document}

% -------- COUVERTURE --------
\begin{titlepage}
\pagecolor{verdescuro}
\color{white}
\vspace*{2.5cm}
\begin{center}
  {\fontsize{36}{42}\selectfont\bfseries Manger Consciemment}\\[1.2cm]
  \textcolor{verdecharo}{\rule{6cm}{0.6pt}}\\[1.2cm]
  {\large\itshape Guide pratique sur les additifs alimentaires}\\[0.6cm]
  {\normalsize Avec E-Codes Reader pour Android}\\[3cm]
  {\normalsize\textcolor{verdecharo}{faustobe.it}}
\end{center}
\vfill
\begin{center}
  {\small Version 1.0\ \textbullet\ 2026\ \textbullet\ CC BY-NC 4.0}
\end{center}
\end{titlepage}
\nopagecolor

\frontmatter
\tableofcontents
\mainmatter

% ============================================================
\chapter{Comment utiliser l'app pour faire ses courses en conscience}

Apprenez à reconnaître les aliments les plus transformés et à choisir ceux
vraiment simples et naturels en quelques minutes, sans devenir un expert
des additifs.

\section{Avant de scanner, faites trois vérifications rapides}

Avant même d'ouvrir l'app, observez l'emballage. Les étiquettes recto sont
souvent conçues pour faire paraître les produits plus sains qu'ils ne le sont.

\begin{itemize}
  \item \textbf{Regardez le recto de l'emballage.} Méfiez-vous des allégations
    comme «sans sucre mais avec édulcorants», «sans matières grasses mais avec
    émulsifiants», «goût extra» ou «aromatisé». Elles indiquent souvent que le
    produit a été fortement modifié.
  \item \textbf{Ouvrez l'app et scannez le produit.} Utilisez l'appareil photo
    pour scanner le code-barres ou importez une photo de l'étiquette. L'app vous
    affichera immédiatement la liste des ingrédients et des additifs présents.
  \item \textbf{Vérifiez trois points clés.} Combien d'additifs apparaissent
    (surtout les codes~E) ? La liste des ingrédients est-elle longue et
    complexe ? Quelle quantité de sucres, sel et graisses est déclarée ?
\end{itemize}

Si vous répondez «beaucoup» ou «longue/compliquée» à plus d'une question, vous
êtes probablement face à un aliment très transformé.

\section{Trois règles pour repérer un aliment «trop transformé»}

\textbf{Peu d'ingrédients = mieux.} Si la liste est courte et contient des noms
familiers (ex. «farine, eau, sel, huile»), le produit est peu transformé et
généralement plus fiable.

\textbf{Beaucoup d'ingrédients «imprononçables».} Des termes comme «amidons
modifiés», «protéines isolées», «mono- et diglycérides d'acides gras»,
«exhausteurs de goût» ou «arômes complexes» signalent clairement un aliment
très transformé.

\textbf{Beaucoup de codes~E et d'additifs.} Tous les~E ne sont pas dangereux,
mais leur présence en masse --- colorants, conservateurs, stabilisants,
émulsifiants --- indique un degré élevé de transformation industrielle.

\section{Comment interpréter le résultat de l'app}

Chaque produit reçoit une évaluation rapide : pensez à ces couleurs comme à
un feu tricolore alimentaire.

\boxvert{Peu transformé --- feu vert}{Peu d'ingrédients, peu ou pas d'additifs~E.
Sucres et sel modérés. Excellent choix pour les courses du quotidien.}

\boxjaune{Modérément transformé --- avec modération}{Quelques~E, arômes naturels
ou sucres ajoutés. Acceptable de temps en temps, pas comme base de chaque repas.}

\boxrouge{Ultra-transformé --- à limiter ou éviter}{Beaucoup d'additifs, sucres
ajoutés, arômes complexes et graisses d'origine industrielle. À réserver comme
vraie exception, surtout pour les enfants.}

\section{Comment utiliser l'app concrètement en faisant ses courses}

\textbf{Étape 1 : Commencez par le rayon frais et petit-déjeuner.} Scannez
yaourts, pain, lait, boissons. Cherchez des produits avec peu d'ingrédients,
peu de codes~E, et des sucres simples (fruits, lait, sucre de canne) plutôt que
des sirops et édulcorants.

\textbf{Étape 2 : Passez ensuite aux sucreries et snacks.} Si la liste est très
longue et pleine de~E, l'app affichera du rouge. Choisissez toujours l'alternative
plus simple : fruits, pain complet, noix.

\textbf{Étape 3 : Utilisez l'app pour comparer des produits similaires.} Souvent
celui qui a moins d'ingrédients est juste légèrement moins attrayant en rayon,
mais bien plus naturel.

\textbf{Étape 4 : Réservez les «rouges» à l'exception, pas à la règle.} Il suffit
de réduire les ultra-transformés. Fixez-vous un seuil --- par exemple, maximum
1--2 produits «rouges» par semaine.

\section{Tableau récapitulatif}

{\small
\begin{tabularx}{\textwidth}{|L|L|}
\hline
\rowcolor{verdescuro}\textcolor{white}{\textbf{Résultat dans l'app}}
  & \textcolor{white}{\textbf{Que faire en pratique}} \\
\hline
Peu d'ingrédients, peu/pas de~E
  & Base quotidienne : yaourt nature, pain complet, coulis de tomates \\
\hline
Quelques~E, sucres modérés
  & Consommez-le de temps en temps, pas comme aliment principal \\
\hline
Beaucoup de~E, sucres ajoutés, arômes complexes
  & Limitez ou évitez : réservez-le à une vraie exception \\
\hline
\end{tabularx}
}

\section{Pourquoi apprendre à lire les additifs fait vraiment la différence}

Reconnaître les aliments les plus transformés n'est pas une «phobie des additifs» :
c'est porter attention à la qualité de ce que nous mangeons. De nombreuses études
suggèrent que réduire les ultra-transformés aide à long terme à contrôler le poids,
à préserver la santé cardiovasculaire et à stabiliser le métabolisme.

% ============================================================
\chapter{Comment lire une étiquette nutritionnelle}

Ingrédients, calories, sel caché : apprenez à déchiffrer une étiquette en moins
d'une minute et arrêtez d'acheter des produits qui semblent sains mais ne le sont pas.

\section{La liste des ingrédients : par où commencer}

Les ingrédients suivent un \textbf{ordre décroissant de poids} : ce qui apparaît
en premier est présent en plus grande quantité.

\begin{itemize}
  \item Le premier ingrédient est le plus important.
  \item Si le premier ingrédient est «sucre», «sirop de glucose» ou «huile de
    palme», le produit est construit autour de cet élément.
  \item La longueur de la liste indique le degré de transformation : pain maison
    = 4 ingrédients ; pain industriel = 15 ou plus.
  \item Cherchez des noms difficiles à prononcer comme «amidon modifié» ou
    «hydrogéné».
\end{itemize}

\boxinfo{Règle pratique}{Si tous les ingrédients ne se trouvent pas dans une
épicerie ordinaire, le produit est probablement très transformé.}

\section{Le tableau nutritionnel : ce qu'il faut regarder}

\begin{description}
  \item[Sucres] Inférieur à 5\,g/100\,g = faible ; entre 5 et 22,5\,g = moyen ;
    au-dessus de 22,5\,g = élevé. Pour les boissons : seuil faible = 2,5\,g/100\,ml.
  \item[Sel] En dessous de 0,3\,g/100\,g = faible ; au-dessus de 1,5\,g/100\,g
    = élevé. Le sodium indiqué sur l'étiquette doit être multiplié par 2,5 pour
    obtenir la valeur en sel.
  \item[Graisses saturées] Au-dessus de 5\,g/100\,g = élevé.
\end{description}

\section{Portions vs 100\,g : l'astuce qui prête à confusion}

Les entreprises mettent en avant les valeurs par portion, souvent sous-estimées.
Utilisez toujours les valeurs pour 100\,g pour comparer entre marques.

\boxjaune{Exemple}{Si l'emballage indique 10\,g de sucres par portion de 30\,g
et que vous en mangez 60\,g, vous consommez 20\,g de sucres.}

\section{Le sel caché : où se cache-t-il vraiment}

Le sel se trouve aussi dans les produits apparemment sucrés.

\begin{itemize}
  \item \textbf{Pain et produits de boulangerie :} deux tranches de pain industriel
    peuvent contenir 0,5--0,8\,g de sel.
  \item \textbf{Céréales du matin :} même les versions «saines» et «complètes»
    peuvent contenir 0,8--1,2\,g de sel pour 100\,g.
  \item \textbf{Sauces et condiments :} entre 1 et 3\,g de sel pour 100\,g.
\end{itemize}

\section{Les allégations nutritionnelles : ce qu'elles signifient vraiment}

\textbf{«Sans sucres ajoutés» :} ne signifie pas sans sucres (peut contenir
fructose, lactose), ni sans édulcorants artificiels.

\textbf{«Light» ou «allégé» :} réduit de 30\,\% par rapport au produit de
référence ; la graisse retirée est souvent remplacée par des sucres ou des
épaississants.

\textbf{«Riche en fibres» ou «source de protéines» :} requiert des seuils
minimaux légaux, mais ne dit rien sur les autres ingrédients.

\section{Liste de contrôle rapide à utiliser au supermarché}

\begin{enumerate}
  \item Regarder le premier ingrédient.
  \item Compter les ingrédients : moins de 5 = excellent ; 5--10 = acceptable ;
    plus de 10--12 = vérifier attentivement.
  \item Vérifier les codes~E (additifs).
  \item Vérifier sel et sucres pour 100\,g.
  \item Ignorer les allégations en grand sur la face avant.
\end{enumerate}

\boxvert{L'astuce la plus simple}{Si vous reconnaissez tous les ingrédients et
pourriez les cuisiner à la maison, vous êtes sur la bonne voie.}

% ============================================================
\chapter{Additifs et ultra-transformés : quelle différence ?}

Un produit avec beaucoup d'additifs est-il nécessairement mauvais ? Et un
ultra-transformé est-il toujours plein d'additifs ? La réponse pourrait vous
surprendre.

\section{Qu'est-ce que les additifs alimentaires}

Les additifs sont des substances intentionnellement incorporées aux aliments pour
assurer une fonction technologique : conservation, coloration, épaississement,
émulsification ou stabilisation. En Europe, chaque additif approuvé reçoit un
code~E et doit passer une évaluation de sécurité de l'EFSA.

\begin{itemize}
  \item Les additifs ne sont pas tous synthétiques : l'acide citrique, la
    curcumine, la lécithine de soja et la vitamine~C sont d'origine naturelle.
  \item La majorité des~E approuvés est considérée comme sûre aux doses
    autorisées.
  \item La problématique réside dans la quantité et la combinaison de nombreux
    additifs.
\end{itemize}

\section{Qu'est-ce que les aliments ultra-transformés}

Le terme «ultra-transformé» vient de la classification NOVA développée par
l'Université de São Paulo. Un produit NOVA~4 contient des ingrédients
introuvables en cuisine domestique : protéines isolées, amidons modifiés,
hydrolysats de protéines, sirops de glucose-fructose, arômes «nature-identiques»,
émulsifiants, colorants et stabilisants.

\textbf{Ultra-transformés :} gâteaux industriels, snacks emballés, céréales
sucrées, boissons gazeuses, saucisses reconstituées, plats surgelés avec
sauces, yaourts aromatisés, pain de mie industriel.

\textbf{Ce qui n'est PAS ultra-transformé :} fromages affinés, charcuteries
de qualité, pain artisanal, pâtes sèches traditionnelles, conserves de tomates
sans additifs, vinaigre, huile d'olive extra vierge.

\section{Quand ils se recoupent --- et quand non}

\begin{description}
  \item[Additifs mais PAS ultra-transformé :] fromage avec conservateur E252,
    vin avec sulfites E220, huile avec vitamine~E E306.
  \item[Ultra-transformé avec PEU d'additifs~E :] pain de mie avec amidon
    modifié et arômes, barres protéinées, surimi.
  \item[Ultra-transformé ET avec de nombreux additifs :] gâteaux industriels,
    sauces industrielles, boissons énergisantes.
\end{description}

\boxinfo{Conclusion}{Un produit avec un~E naturel n'est pas un problème. Un
produit avec 15 ingrédients industriels sans aucun~E est quand même
un ultra-transformé.}

\section{Le double problème : pourquoi les deux comptent}

\textbf{Risques liés à des additifs spécifiques :} colorants azoïques et
hyperactivité infantile, nitrates et nitrosamines cancérigènes, édulcorants
intenses et altération du microbiome.

\textbf{Risques liés à l'ultra-transformation :} des études à grande échelle
(NutriNet-Santé, UK Biobank) associent une consommation élevée de NOVA~4 à
un risque accru d'obésité, de diabète de type~2, de maladies cardiovasculaires
et de certains cancers.

\textbf{L'effet de combinaison :} dans la vie réelle, nous consommons des
combinaisons de 10 à 20 additifs différents au cours du même repas.

\section{Comment les reconnaître en pratique}

\begin{enumerate}
  \item Repérez les ingrédients «de laboratoire» : amidons modifiés, protéines
    isolées, hydrolysats, sirops de glucose-fructose.
  \item Comptez les codes~E et identifiez les catégories critiques.
  \item Tenez compte du contexte du repas.
  \item Utilisez l'application comme aide, pas comme juge absolu.
\end{enumerate}

% ============================================================
\chapter{Classification NOVA : qu'est-ce que c'est et comment l'utiliser}

NOVA ne mesure pas les nutriments mais le \emph{degré de transformation
industrielle}. Comprendre les 4 groupes change la façon dont vous faites
vos courses.

\section{Qu'est-ce que la classification NOVA}

NOVA est un système de classification développé par le professeur Carlos Monteiro
à l'Université de São Paulo et adopté par l'OMS et la FAO. Contrairement au
Nutri-Score, NOVA classe selon le \emph{degré de transformation industrielle}
et non selon la composition nutritionnelle.

Les procédés modernes peuvent altérer les fibres, la biodisponibilité des
nutriments et les réponses hormonales, indépendamment de la composition finale.
NOVA divise les aliments en 4 groupes distincts, non en une échelle continue.

\section{Groupe 1 --- Aliments non transformés ou peu transformés}

Aliments sans transformation industrielle significative, ou seulement avec des
procédés physiques minimaux : séchage, mouture, réfrigération, pasteurisation.

\textbf{Exemples :} fruits et légumes frais ou surgelés, viandes et poissons
frais sans additifs, œufs, légumineuses sèches, céréales en grain (riz, avoine,
épeautre), farine complète, lait frais pasteurisé, yaourt nature sans additifs,
thé, café, eau.

\boxvert{Objectif}{Ces aliments devraient constituer la majorité de chaque repas.}

\section{Groupe 2 --- Ingrédients culinaires transformés}

Substances extraites d'aliments du Groupe~1 utilisées en cuisine pour cuire,
assaisonner et préparer les aliments ; jamais consommées seules.

\textbf{Exemples :} huile d'olive, beurre, vinaigre, sel, sucre, miel, farine
blanche, amidon de maïs pur, sirop d'érable.

À utiliser avec modération : le sucre, le sel et les graisses saturées en
excès restent défavorables.

\section{Groupe 3 --- Aliments transformés}

Produits résultant de l'ajout d'ingrédients du Groupe~2 à des aliments du
Groupe~1, avec des procédés comme la conservation au sel, le fumage, la
fermentation ou l'affinage. Ils contiennent généralement 2--3 ingrédients.

\textbf{Exemples :} légumes et légumineuses en conserve (avec eau et sel),
fruits confits, noix grillées et salées, anchois à l'huile, jambon cru de
qualité, fromages traditionnels (parmesan, comté), pain artisanal, vin, bière
artisanale.

Le Groupe~3 n'est pas à éviter ; de nombreux aliments traditionnels
y appartiennent.

\section{Groupe 4 --- Ultra-transformés : le groupe à limiter}

Formulations industrielles ne contenant pas d'aliments entiers ou n'en utilisant
qu'une infime partie. Ils sont construits avec des ingrédients extraits, purifiés
ou modifiés chimiquement, plus des additifs technologiques en grandes quantités.

Signal clé : ils contiennent des \textbf{ingrédients qu'on ne trouve pas dans
une cuisine domestique}.

Ingrédients typiques : protéines hydrolysées, amidons modifiés, sirop de
glucose-fructose, huiles hydrogénées, arômes «nature-identiques», colorants
artificiels, édulcorants intenses (aspartame, saccharine, sucralose),
émulsifiants (E471, E472).

\boxrouge{Données clés}{En France, environ 30--35\,\% des calories journalières
proviennent d'ultra-transformés. Une consommation supérieure à 20\,\% des
calories journalières de NOVA~4 est associée à des risques pour la santé.}

\section{NOVA dans les courses au quotidien}

\textbf{1.~«Pourrais-je le cuisiner à la maison avec des ingrédients normaux ?»}
Si non, c'est probablement du Groupe~4.

\textbf{2.~Cherchez les «mots révélateurs» du Groupe~4 :} «amidon modifié»,
«protéines de soja isolées», «maltodextrines», «sirop de\ldots», «arômes
complexes», «huile de palme hydrogénée».

\textbf{3.~Règle pratique :} plus de la moitié de l'assiette doit provenir
du Groupe~1.

\section{Tableau récapitulatif NOVA}

{\small
\begin{tabularx}{\textwidth}{|l|L|L|l|}
\hline
\rowcolor{verdescuro}
  \textcolor{white}{\textbf{Groupe}}
  & \textcolor{white}{\textbf{Définition}}
  & \textcolor{white}{\textbf{Exemples}}
  & \textcolor{white}{\textbf{Alimentation}} \\
\hline
NOVA~1 & Non/peu transformés & Fruits, légumes, œufs, viande fraîche & Base \\
\hline
NOVA~2 & Ingrédients culinaires & Huile, sel, sucre, farine, beurre & Modération \\
\hline
NOVA~3 & Transformés traditionnels & Conserves, fromages, jambon, pain artisanal & Acceptable \\
\hline
NOVA~4 & Ultra-transformés & Gâteaux industriels, saucisses, céréales sucrées & À limiter \\
\hline
\end{tabularx}
}

% ============================================================
\chapter{Ingrédients à surveiller sur les étiquettes}

Pas besoin de lire chaque étiquette comme un chimiste. Quelques termes clés
suffisent à éviter les produits les plus problématiques et à faire de meilleurs
choix en quelques secondes.

\section{Sucres : les 10 noms cachés sur les étiquettes}

Le sucre utilise de nombreux alias pour sembler moins important dans la liste
des ingrédients. Les 10 alias les plus courants :

\begin{itemize}
  \item Sirop de glucose-fructose
  \item Maltodextrines
  \item Dextrose
  \item Fructose
  \item Jus de canne concentré
  \item Sirop de maïs / sirop de riz
  \item Miel déshydraté
  \item Sucre de coco
  \item Nectar d'agave
\end{itemize}

\boxjaune{Le piège du «sans sucres ajoutés»}{Un jus «100\,\% fruit sans sucres
ajoutés» peut contenir 10--12\,g de sucres pour 100\,ml --- plus qu'un Coca-Cola
(10,6\,g/100\,ml). L'OMS recommande de maintenir les sucres libres sous 10\,\%
des calories journalières (environ 50\,g/jour), de préférence sous 5\,\% (25\,g).}

\section{Graisses à éviter : hydrogénées et trans}

\textbf{Graisses hydrogénées et partiellement hydrogénées :} recherchez «huile
végétale partiellement hydrogénée», «graisse végétale hydrogénée», «shortening».
En Europe, ils sont interdits au-dessus de 2\,g/100\,g de graisses depuis 2021.

\textbf{Huile de palme :} saturée mais non hydrogénée. Lors du raffinage à
hautes températures, des composés potentiellement cancérigènes se forment
(3-MCPD, glycidol).

\boxvert{Les graisses qui conviennent}{Huile d'olive extra vierge, huile de lin,
noix et oléagineux, avocat, graisses des poissons gras.}

\section{Additifs controversés : les codes E à connaître}

\textbf{Colorants azoïques --- vigilance maximale avec les enfants :}
E102 (Tartrazine), E104 (Jaune de quinoléine), E110 (Jaune orangé~S),
E122 (Carmoisine), E124 (Rouge Ponceau 4R), E129 (Rouge Allura~AC).
Étude EFSA 2007 : corrélation avec l'hyperactivité chez les enfants. Mention
obligatoire sur l'étiquette.

\textbf{Nitrates et nitrites (E249--E252) :} utilisés dans les charcuteries,
saucisses, bacon. En présence de protéines animales, ils forment des
nitrosamines, classées comme probables cancérigènes (IARC groupe 2A).

\textbf{Édulcorants intenses (E950, E951, E952, E954, E955) :} sûrs aux doses
approuvées, mais l'OMS (2023) a exprimé une mise en garde sur l'usage
chronique pour le contrôle du poids.

\section{Sodium caché : au-delà du sel de cuisine}

Environ 75\,\% du sodium dans l'alimentation occidentale n'est pas ajouté
pendant la cuisson mais est déjà présent dans les produits emballés.

Principales sources : pain (première source), fromages, charcuteries, plats
préparés, sauces en bocal, bouillon cube, sauce soja, ketchup, mayonnaise,
snacks salés, céréales industrielles.

\boxinfo{Conversion sodium--sel}{Multipliez la valeur de sodium par 2,5.
Exemple : 0,6\,g de sodium = 1,5\,g de sel. Limite OMS : 5\,g de sel par
jour (environ 2\,g de sodium).}

\section{La règle des 5 : le filtre le plus rapide}

\begin{enumerate}
  \item \textbf{Plus de 5 ingrédients ?} Commencez à lire attentivement.
  \item \textbf{Sucres supérieurs à 5\,g/100\,g ?} Vérifiez s'il est le premier
    ou deuxième ingrédient.
  \item \textbf{Sel supérieur à 0,5\,g/100\,g dans un produit sucré ?} Signal
    d'ultra-transformation.
  \item \textbf{Plus de 5 additifs~E ?} Le produit est probablement ultra-transformé.
  \item \textbf{Vous ne reconnaissez pas 5 ingrédients ou plus ?} Reposez-le
    sur l'étagère.
\end{enumerate}

\boxinfo{Note}{La règle des 5 est une orientation, non un verdict absolu.}

% ============================================================
\chapter{Aliments recommandés : les options les plus saines au supermarché}

Il ne s'agit pas d'être parfait ni de ne manger que de la salade. Il s'agit
de savoir quels produits, à commodité égale, sont vraiment meilleurs que
les autres.

\section{Le principe de base : substituer, pas éliminer}

\begin{description}
  \item[Même temps, moins de transformation.] Remplacer yaourt aromatisé par
    yaourt nature, riz précuit par riz complet, plats préparés par légumineuses
    en boîte.
  \item[Règle 80/20.] Si 80\,\% de l'alimentation est peu transformée, les
    20\,\% restants ne posent pas problème.
  \item[Accessibilité tarifaire.] Pois chiches en boîte, œufs, farine complète,
    lentilles coûtent moins que les produits transformés équivalents.
\end{description}

\section{Céréales et glucides : les meilleurs choix}

\textbf{Recommandées :} flocons d'avoine nature, riz complet, épeautre, orge
perlé, sarrasin, pâtes complètes, pain avec maximum 4--5 ingrédients.

\textbf{À limiter :} céréales sucrées, pain de mie avec additifs, crackers à
l'huile de palme, riz précuit avec additifs, bollería industrielle.

\boxjaune{Astuce}{Si les sucres dépassent 10\,g/100\,g, vous mangez un produit
de dessert, pas un petit-déjeuner sain.}

\section{Protéines : les meilleures sources}

\textbf{Poissons et fruits de mer :} frais ou surgelés sans additifs, sardines
et maquereaux à l'huile d'olive, thon au naturel, anchois en sel.

\textbf{Œufs et produits laitiers :} œufs entiers, yaourt grec nature, fromage
blanc, fromages affinés modérément. Évitez les yaourts «protéinés» avec
épaississants et édulcorants.

\textbf{Légumineuses :} lentilles, pois chiches, haricots en boîte au naturel,
riches en fibres et acides aminés.

\section{Bonnes graisses : les plus simples à ajouter}

\textbf{Recommandées :} huile d'olive extra vierge, noix, amandes, graines de
lin, avocat, poissons gras.

\textbf{Avec modération :} beurre naturel, huile de coco, lait entier.

\textbf{À éviter :} margarines hydrogénées, shortening, huile de palme
raffinée, crème végétale avec émulsifiants.

\section{Encas intelligents}

\begin{enumerate}
  \item \textbf{Fruits secs nature :} noix, amandes, noisettes sans sel
    (30\,g apportent protéines et minéraux).
  \item \textbf{Fruits frais entiers :} pomme, poire, kiwi (contiennent des
    fibres contrairement aux jus).
  \item \textbf{Yaourt grec nature :} haute protéine, faible indice glycémique.
  \item \textbf{Fromage à pâte dure + crackers complets :} satiété durable
    avec ingrédients simples.
\end{enumerate}

\section{Liste de courses type}

{\small
\begin{tabularx}{\textwidth}{|l|L|}
\hline
\rowcolor{verdescuro}
  \textcolor{white}{\textbf{Catégorie}}
  & \textcolor{white}{\textbf{Produits}} \\
\hline
Céréales & Flocons d'avoine, riz complet, pâtes complètes, pain artisanal \\
\hline
Légumineuses & Lentilles, pois chiches, haricots \\
\hline
Protéines & Œufs, sardines, thon, yaourt grec nature \\
\hline
Légumes & Frais ou surgelés sans additifs \\
\hline
Graisses & Huile EVO, noix, amandes, graines \\
\hline
Fruits & Fruits entiers de saison, baies surgelées \\
\hline
Condiments & Vinaigre, sauce soja peu transformée, épices \\
\hline
\end{tabularx}
}

\boxvert{Test du caddie}{Vérifier que plus de la moitié des produits contient
moins de 5 ingrédients.}

% ============================================================
\chapter{Pourquoi les aliments ultra-transformés sont mauvais pour la santé}

Ce n'est pas seulement une question de calories ou de graisses : les
ultra-transformés perturbent la satiété, altèrent le microbiome et déclenchent
une inflammation chronique.

\section{Obésité et altération du métabolisme}

Les UPF sont conçus pour être \emph{hyper-appétissants} et interfèrent
activement avec les signaux naturels de satiété.

\boxinfo{Étude Kevin Hall (NIH, 2019)}{Les participants au régime
ultra-transformé consommaient \emph{en moyenne 500 calories de plus par jour}
que ceux au régime d'aliments peu transformés, avec gain de près d'un kilogramme
en deux semaines.}

Les UPF contiennent en moyenne 2,15\,kcal par gramme (presque le double des
aliments frais) et réduisent le peptide~YY, l'hormone intestinale qui signale
la satiété au cerveau.

Statistique : \textbf{+41\,\% de risque d'obésité abdominale} pour les
consommateurs réguliers d'UPF.

\section{Inflammations et dysfonctionnements immunitaires}

La consommation fréquente d'ultra-transformés est associée à une augmentation
des réponses immunitaires inflammatoires, au-delà du profil nutritionnel seul.

Des additifs fréquents dans les UPF --- carboxyméthylcellulose (E466),
polysorbate~80 (E433), dioxyde de titane (E171) --- peuvent altérer la muqueuse
intestinale et augmenter la perméabilité intestinale (\textit{leaky gut}).

\textbf{Maladies auto-immunes associées :} maladie cœliaque, thyroïdite de
Hashimoto, sclérose en plaques, lupus érythémateux systémique, diabète
de type~1.

\section{Santé intestinale et microbiote}

Le processus d'ultra-transformation prive les aliments de fibres alimentaires
et de composés végétaux bioactifs. Les émulsifiants et les édulcorants
artificiels affaiblissent la couche muqueuse qui tapisse l'intestin.

\textbf{Pathologies fréquentes :} maladie de Crohn, rectocolite hémorragique,
syndrome de l'intestin irritable, ulcères gastriques, polypes précancéreux
du côlon.

\section{Un risque systémique : plusieurs mécanismes à la fois}

{\small
\begin{tabularx}{\textwidth}{|l|L|L|}
\hline
\rowcolor{verdescuro}
  \textcolor{white}{\textbf{Mécanisme}}
  & \textcolor{white}{\textbf{Effet}}
  & \textcolor{white}{\textbf{Conséquence}} \\
\hline
Hyper-appétibilité & Consommation excessive & Obésité, diabète type~2 \\
\hline
Carence en fibres & Dysbiose & Inflammation chronique \\
\hline
Additifs (émulsifiants) & Perméabilité intestinale & MICI, maladies auto-immunes \\
\hline
Sucres/édulcorants & Résistance insulinique & Diabète, syndrome métabolique \\
\hline
Remplacement des frais & Carence micronutriments & Stress oxydatif \\
\hline
\end{tabularx}
}

Les grandes études --- NutriNet-Santé (France, 100\,000+ participants) et
UK Biobank (500\,000+) --- associent le consommation élevée d'UPF à des
risques accrus de cancer, maladies cardiovasculaires, dépression et mortalité.

\section{Que faire en pratique}

Le message pratique : \emph{réduisez la part d'UPF} plutôt que d'éliminer
complètement.

\begin{enumerate}
  \item \textbf{Identifiez vos UPF habituels :} céréales sucrées, snacks,
    pain de mie, plats préparés, boissons aromatisées.
  \item \textbf{Remplacez-les progressivement :} avoine pour céréales, yaourt
    nature pour yaourt aromatisé, pain artisanal.
  \item \textbf{Augmentez les fibres --- priorité :} légumineuses, légumes,
    céréales complètes, fruits entiers.
  \item \textbf{Utilisez l'app :} E-Codes Reader identifie les ultra-transformés
    lors du scan du code-barres.
  \item \textbf{Ne visez pas la perfection :} même une amélioration de
    20--30\,\% produit des effets mesurables.
\end{enumerate}

% ============================================================
\chapter{Sécurité alimentaire pour les enfants et les personnes âgées}

Les enfants et les personnes âgées réagissent différemment aux additifs et
ingrédients de mauvaise qualité par rapport aux adultes en bonne santé.

\section{Pourquoi les enfants et les personnes âgées sont plus vulnérables}

\textbf{Enfants : poids corporel faible, systèmes immatures.} Un enfant de 15\,kg
consommant un goûter de 30\,g avec colorants reçoit une dose relative 4 à 5 fois
supérieure à celle d'un adulte de 70\,kg. Le foie et les reins métabolisent
certaines substances plus lentement chez les jeunes enfants.

\textbf{Personnes âgées : fonctions rénale et hépatique réduites.} La capacité
d'éliminer certaines substances diminue avec l'âge. La polymédication peut
interagir avec des additifs. Le sodium excessif aggrave l'hypertension et
l'insuffisance rénale.

\textbf{Effets cumulatifs :} une consommation quotidienne de produits contenant
certains colorants produit des effets cumulatifs significatifs.

\section{Additifs critiques pour les enfants}

\textbf{Colorants azoïques --- à éviter :}
E102, E104, E110, E122, E124, E129 nécessitent l'étiquetage : «peut avoir des
effets indésirables sur l'activité et l'attention des enfants». Retrouvés dans
les bonbons, boissons oranges/rouges, gelées, céréales du matin.

\textbf{Édulcorants intenses --- non approuvés pour jeunes enfants :}
E951, E950, E955, E954. Présents dans les yaourts «0\,\% sucre», boissons
\textit{light}, sirops, bonbons sans sucre.

\boxrouge{Caféine}{L'EFSA recommande maximum 3\,mg/kg/jour. Une canette de
250\,ml d'energy drink peut contenir jusqu'à 80\,mg --- presque le double
de la limite pour un enfant de 15\,kg.}

\section{Ce qu'il faut limiter chez les personnes âgées}

\textbf{Sodium : le risque principal.} Objectif : moins de 5\,g de sel
quotidien. Vecteurs principaux : plats préparés, charcuteries, fromages,
bouillons industriels.

\textbf{Phosphates (E338--E452) :} courants dans les viandes transformées,
fromages fondus. Une consommation élevée interfère avec l'absorption du calcium
--- un problème pour les personnes à risque d'ostéoporose.

\textbf{Interactions médicamenteuses :} la vitamine~K des légumes verts avec
les warfarines ; le pamplemousse avec les statines ; le jus de cranberry avec
certains anticoagulants.

\section{Produits «pour enfants» à regarder attentivement}

\textbf{Céréales du matin «pour enfants» :} contiennent souvent 25--35\,g de
sucres pour 100\,g, colorants artificiels et vitamines ajoutées pour
compenser la faible qualité nutritive.

\textbf{Boissons «aux jus de fruits» :} peuvent ne contenir que 10--30\,\% de
jus réel. Vérifiez toujours le pourcentage de fruits sur l'étiquette.

\textbf{Yaourts «pour enfants» aux fruits :} contiennent souvent plus de sucre
qu'un yaourt standard, plus d'épaississants et arômes artificiels.

\section{Liste sécurisée pour les courses}

\begin{enumerate}
  \item \textbf{Pour les enfants :} produits avec peu d'ingrédients reconnaissables ;
    évitez plus de 2--3 colorants ou édulcorants artificiels.
  \item \textbf{Vérifiez les sucres dans les produits «pour enfants» :}
    moins de 10\,g/100\,g.
  \item \textbf{Pour personnes âgées :} priorité à faible sodium ;
    moins de 0,3\,g sel/100\,g.
  \item \textbf{Scannez avec l'app avant d'acheter :} E-Codes Reader indique
    immédiatement les colorants azoïques et le niveau de sodium.
  \item \textbf{Privilégiez le Groupe~1--2 (NOVA) pour les repas principaux :}
    légumes, légumineuses, œufs, poisson, yaourt nature, fruits entiers.
\end{enumerate}

\section{Tableau récapitulatif}

{\small
\begin{tabularx}{\textwidth}{|L|l|l|}
\hline
\rowcolor{verdescuro}
  \textcolor{white}{\textbf{Additif/ingrédient}}
  & \textcolor{white}{\textbf{Enfants}}
  & \textcolor{white}{\textbf{Personnes âgées}} \\
\hline
Colorants azoïques (E102, E110, E122\ldots) & À éviter & À limiter \\
\hline
Édulcorants intenses & Pas pour jeunes enf. & Avec modération \\
\hline
Sodium élevé (>1,5\,g sel/100\,g) & À limiter & À éviter \\
\hline
Phosphates (E338--E452) & Modération & À limiter \\
\hline
Caféine & À éviter <12 ans & Attention \\
\hline
Nitrates/nitrites (E249--E252) & À limiter & À limiter \\
\hline
\end{tabularx}
}

% ============================================================
\backmatter

\chapter*{Crédits et téléchargement}
\addcontentsline{toc}{chapter}{Crédits}

\begin{center}
\vspace{1cm}
{\large\bfseries\color{verdescuro} Manger Consciemment}\\[0.4cm]
{\itshape Guide pratique sur les additifs alimentaires}\\[0.3cm]
{\small Avec E-Codes Reader pour Android}

\vspace{1.5cm}
\textbf{Version 1.0} \quad \textbullet \quad 2026\\
\textbf{Licence :} Creative Commons BY-NC 4.0\\[0.3cm]
\textbf{Site web :} \href{https://faustobe.it}{faustobe.it}\\
\textbf{Application :} \href{https://play.google.com/store/apps/details?id=it.faustobe.ecodes}{E-Codes Reader sur Google Play}

\vspace{1.5cm}
{\small Scannez le code QR pour télécharger l'application :}\\[0.5cm]

\qrcode[height=4cm]{https://play.google.com/store/apps/details?id=it.faustobe.ecodes}

\vspace{1cm}
{\footnotesize Cet ebook a été produit en \LaTeX{} par faustobe.}
\end{center}

\end{document}
